% ===== CipherFlow Progress Report =====
\documentclass[a4paper,12pt]{report}

\usepackage[french]{babel}
\usepackage[T1]{fontenc}
\usepackage[utf8]{inputenc}

\usepackage{geometry}
\geometry{left=2.5cm,right=2.5cm,top=2cm,bottom=3.5cm}

\usepackage{graphicx}
\usepackage{float}
\usepackage{hyperref}
\usepackage{wallpaper}

\usepackage{booktabs}
\usepackage[table]{xcolor}
\usepackage{colortbl}

\usepackage{array}
\usepackage{tabularx}
\usepackage{longtable}
\usepackage{pdflscape}
\usepackage{enumitem}
\usepackage{listings}
\usepackage{tcolorbox}
\tcbuselibrary{skins, breakable}

\definecolor{FSTblue}{RGB}{0,102,204}
\definecolor{codebg}{RGB}{245,245,245}
\definecolor{codeframe}{RGB}{200,200,200}

\lstdefinelanguage{Kotlin}{
  sensitive=true,
  keywords={fun, val, var, when, do, while, if, else, for, in, return,
            true, false, null, class, data, object, companion, override,
            import, package, is, as, by, typealias, this, super, init,
            constructor, abstract, open, sealed, interface, enum, break,
            continue, throw, try, catch, finally, it},
  morestring=[b]",
  morestring=[b]',
  morecomment=[l]{//},
  morecomment=[s]{/*}{*/},
}

\lstdefinestyle{CodeStyle}{
  basicstyle=\ttfamily\small,
  backgroundcolor=\color{codebg},
  frame=single,
  rulecolor=\color{codeframe},
  breaklines=true,
  showstringspaces=false,
  keywordstyle=\color{FSTblue}\bfseries,
  stringstyle=\color{teal},
  commentstyle=\color{gray}\itshape,
  numbers=left,
  numberstyle=\tiny\color{gray},
  numbersep=5pt,
  xleftmargin=10pt,
}

\frenchbsetup{StandardLists=true}

\newcommand{\question}[1]{\section*{#1}\addcontentsline{toc}{section}{#1}}

\newcommand{\screenshot}[3][0.75]{%
  \IfFileExists{./imgs/#2}{%
    \begin{figure}[H]\centering
      \includegraphics[width=#1\textwidth]{./imgs/#2}
      \caption{#3}
    \end{figure}}{%
    \begin{figure}[H]\centering
      \fbox{\parbox[c][5cm][c]{#1\textwidth}{\centering
        \textbf{CAPTURE A INSERER}\\[1ex]\texttt{imgs/\detokenize{#2}}}}
      \caption{#3}
    \end{figure}}%
}

%%%%%%%%%%%%%%%%%%%%%%%%%%%%
% Custom Title Page Setup  %
%%%%%%%%%%%%%%%%%%%%%%%%%%%%
\newcommand*{\makeCustomTitle}{
    \newgeometry{margin=2.5cm}
    \begin{titlepage}
        \ThisLRCornerWallPaper{1.2}{./background.jpg}
        \setlength{\parindent}{0pt}

        \begin{minipage}[t][0.25\textheight]{\textwidth}
            \begin{center}
                \includegraphics[height=0.15\textheight]{./LOGO_FSTM.png}
            \end{center}
            \vspace{2ex}
            \begin{center}
                {
                    \scshape
                    {\LARGE Universit\'{e} Cadi Ayyad\\}
                    {\Large Facult\'{e} des Sciences et Techniques de Marrakech\\[1ex]}
                    {\large D\'{e}partement d'Informatique}
                }
            \end{center}
        \end{minipage}

        \begin{minipage}{\textwidth}
            \vspace{0.15\textheight}
            \begin{center}
                {\bfseries\huge Rapport d'Avancement --- CipherFlow}
            \end{center}
            \vspace{3ex}
        \end{minipage}

        \begin{center}
            {\Large Module : Developpement Frontend}

            \vfill

            \begin{tabular}{@{}p{0.49\textwidth}@{\hspace{0.02\textwidth}}p{0.49\textwidth}@{}}
                \begin{minipage}[t]{0.49\textwidth}
                    \large
                    \textbf{\'{E}tudiants}\\
                    Mohamed \textsc{EL BOUCHIR}\\
                    Bilal \textsc{AKENSOUSS}\\
                    Mohammed \textsc{ALHAYANE}\\
                    Salam \textsc{ELOUZI}
                \end{minipage}
            \end{tabular}

            \vfill

            {\large\today}
        \end{center}
    \end{titlepage}
    \restoregeometry
}

\begin{document}

\makeCustomTitle

\tableofcontents
\newpage

\chapter*{CipherFlow --- Tableau de Bord Anime de Transformation de Donnees}
\addcontentsline{toc}{chapter}{CipherFlow --- Tableau de Bord Anime de Transformation de Donnees}

CipherFlow est une application web de visualisation d'algorithmes cryptographiques
et de transformation de donnees, inspiree de CyberChef. L'objectif principal est
de rendre la composition de pipelines d'operations interactive et educative
grace a des animations pas-a-pas declaratives. Ce rapport decrit la mise en place
de l'interface initiale et du systeme d'extension que chaque membre utilisera pour
ajouter ses propres operations.

\question{Presentation du projet et contexte}

CipherFlow se positionne comme une alternative animee a CyberChef. La difference
fondamentale est la suivante : chaque operation du pipeline possede un composant
\texttt{Visualizer} qui rend, image par image, ce qui se passe a l'interieur de
l'algorithme. L'utilisateur peut ainsi voir, lettre par lettre ou octet par octet,
comment son texte est transforme.

Les technologies utilisees sont :

\begin{itemize}
  \item \textbf{React 18 + TypeScript} --- composants fonctionnels et typage strict.
  \item \textbf{Vite} --- serveur de developpement et bundle de production.
  \item \textbf{Framer Motion (motion/react)} --- animations declaratives fluides.
  \item \textbf{Tailwind CSS} --- mise en page utilitaire sans CSS personnalise.
  \item \textbf{Lucide React} --- icones SVG.
\end{itemize}

\screenshot[0.95]{screenshot1.png}{Interface initiale de CipherFlow au demarrage}

\question{Architecture de l'interface initiale}

L'interface principale est composee de cinq blocs empiles verticalement :

\begin{enumerate}
  \item \textbf{TopBar} --- navigation entre la vue Workbench et la vue Presets,
        bouton Play global.
  \item \textbf{InputPanel} --- zone de saisie du texte d'entree.
  \item \textbf{RecipePipeline} --- pipeline glisser-deposer des operations,
        avec controles de parametres par etape (activer/desactiver, deplacer,
        supprimer), bouton d'ajout, mode dechiffrement, inversion d'ordre,
        sauvegarde de recette.
  \item \textbf{AnimationStage} --- scene d'animation : barre de navigation
        entre etapes, lecteur image par image avec vitesse reglable (0.5x, 1x, 2x),
        rendu du composant \texttt{Visualizer} propre a chaque algorithme.
  \item \textbf{OutputPanel} --- affichage de la sortie finale et de la sortie
        partielle de l'etape active, bouton \og Utiliser comme entree \fg{}.
\end{enumerate}

\screenshot[0.95]{screenshot2.png}{Pipeline avec une operation Caesar active et son animation}

\question{Systeme d'extension : comment ajouter une operation}

Toute la logique de decouverte des algorithmes repose sur un registre central :
\texttt{src/algorithms/registry.ts}. Ajouter une nouvelle operation se fait en
trois etapes sans modifier aucun autre fichier existant :

\begin{enumerate}
  \item Creer un fichier \texttt{src/algorithms/monOperation.tsx} exportant un
        objet de type \texttt{AlgorithmDefinition}.
  \item L'importer dans \texttt{registry.ts} et l'ajouter au tableau
        \texttt{ALGORITHMS}.
  \item Le composant \texttt{Visualizer} inclu dans la definition est
        automatiquement utilise par \texttt{AnimationStage} pour afficher
        l'animation.
\end{enumerate}

L'interface \texttt{AlgorithmDefinition} impose les champs suivants :

\begin{lstlisting}[style=CodeStyle, language={[Sharp]C}, caption={Interface AlgorithmDefinition (TypeScript)}]
export interface AlgorithmDefinition<TState = any> {
  id: string;             // identifiant unique, ex: 'caesar-cipher'
  name: string;           // libelle affiche dans l'UI
  category: AlgorithmCategory;  // 'Cipher' | 'Encoding' | 'Bitwise' | 'Text'
  summary: string;        // description courte
  animationPattern: string;     // tag du style d'animation
  complexity: string;     // ex: 'O(n)'
  paramSchema: ParamField[];    // schema des parametres configurables
  defaultParams: Record<string, any>;
  computeTrace: (input: string, params: Record<string, any>)
    => AlgorithmTrace<TState>;  // logique de calcul + frames
  invertParams?: (params: Record<string, any>) => Record<string, any>;
  Visualizer: ComponentType<VisualizerProps<TState>>;  // composant React
}
\end{lstlisting}

La fonction \texttt{computeTrace} est la cle du systeme : elle prend le texte
d'entree et les parametres, et retourne un objet \texttt{AlgorithmTrace}
contenant un tableau de \texttt{AnimationFrame}. Chaque frame decrit l'etat
intermediaire de la transformation (texte produit jusqu'ici, titre de l'etape,
explication, et un objet \texttt{state} specifique a l'algorithme pour le
\texttt{Visualizer}).

\question{Operations implementees dans la version initiale}

Quatre operations ont ete implementees pour servir de base et de modeles
pour les membres de l'equipe :

\begin{itemize}
  \item \textbf{Caesar Cipher} (\texttt{caesar.tsx}) --- decalage de chaque
        caractere dans l'alphabet. Animation : mise en surbrillance lettre par
        lettre avec affichage du decalage.
  \item \textbf{Base64 Encode/Decode} (\texttt{base64.tsx}) --- encodage Base64
        standard. Animation : groupement des caracteres par blocs de 3 puis
        conversion en 4 caracteres Base64.
  \item \textbf{XOR Cipher} (\texttt{xor.tsx}) --- XOR octet par octet avec
        une cle textuelle. Animation : affichage des valeurs binaires de l'entree,
        de la cle et du resultat XOR.
  \item \textbf{Reverse Text} (\texttt{reverse.tsx}) --- inversion de la chaine
        avec option de permutation de casse. Animation : deplacement visuel des
        caracteres vers leurs positions finales.
\end{itemize}

\screenshot[0.95]{screenshot3.png}{Pipeline multi-etapes avec sortie finale visible}

\question{Repartition du travail et prochaines etapes}

La base commune (interface, registre, composants UI, animations generiques) a
ete construite collectivement. Chaque membre de l'equipe prendra en charge
l'implementation d'au moins une operation supplementaire en suivant le contrat
\texttt{AlgorithmDefinition} :

\begin{itemize}
  \item \textbf{Mohamed El Bouchir} --- operation et \texttt{Visualizer}
        personnalises (a definir).
  \item \textbf{Bilal Akensouss} --- operation et \texttt{Visualizer}
        personnalises (a definir).
  \item \textbf{Mohammed Alhayane} --- operation et \texttt{Visualizer}
        personnalises (a definir).
  \item \textbf{Salam Elouzi} --- operation et \texttt{Visualizer}
        personnalises (a definir).
\end{itemize}

Chaque membre est libre de choisir le style visuel de son \texttt{Visualizer}
(transition par fondu, deplacement, retournement de carte, effet de pulsation,
etc.) du moment que la logique de \texttt{computeTrace} reste pure et testable.

\bigskip

Ce rapport constitue le premier jalon du projet. La prochaine etape sera la
presentation de chaque operation avec son animation lors de la seance suivante.

\end{document}
